% SAMPLE ARTICLE FOR COMBINATORIAL THEORY
%
% CT papers *must* follow this sample file.
% Please, do not erase comments in the head of this sample file as
% they are used to process your file semi-automatically.

% COPYRIGHT NOTICE
% We use the CC BY licence. It allows others to freely distribute 
% your work while giving credit to you. In particular CC BY allows 
% remixing, tweaking and building upon your work.  
% See https://creativecommons.org/licenses/ for a full explanation
% of the Creative Commons licenses.

% DOCUMENT CLASS AND CT.STY
% All CT papers start with these two lines. Do not change make
% any changes in the ct.sty file as we use our own copy.
\documentclass[12pt]{article}
\usepackage{ct}

% GEOMETRY AND SPACES
% To keep the style homogeneous among all CT papers, please
%   * remove all commands that change parameters such as margins or page sizes. 
%   * avoid using any vspace, hspace, newpage, and other spacing commands.

% PACKAGES
% The following packages are already loaded in the style file:
%   * amssymb, amsthm, amsmath for mathematics,
%   * graphicx for images,
%   * hyperref and cleveref for references.
% Load your own packages below only if they are really needed:
% \usepackage{...}

% THEOREM ENVIRONMENTS
% The following environments are already declared in the style file:
%   * theorem-like environments: theorem, lemma, corollary, proposition, fact, observation, claim
%   * definition-like environments: definition, example, conjecture, open, problem, question
%   * remark-like environments: remark, note
% Declare your own environments below only if needed:
% \newtheorem{myenv}[theorem]{MyEnvironment}

% SPECIFICATIONS
% The editors will insert the specifications of your paper.
\specs{n}{vol (issue)}{year}{}

% DATES
% Give the submission and acceptance dates in the format shown.
% The editors will insert the publication date in the third argument.
\dateline{Jan 1, 2020}{Jan 2, 2020}{TBD}

% KEYWORDS
% Give one or more keywords separated by commas. No dot at the end.
\keywords{Broglington manifolds, permutations}

% SUBJECT CODES
% Give one or more subject codes separated by commas. No dot at the end.
% Codes are available from http://www.ams.org/mathscinet/freeTools.html
\MSC{05C88, 05C89}

% TITLE
% If needed, include a line break (\\) at an appropriate place in the title.
\title{An elementary proof\\ of the reconstruction conjecture}

% AUTHORS AND AFFILIATIONS
% Input author, affiliation, address and support information as follows.
% The address should include the country, but does not have to include the street address. 
% Give at least one email address.
% Grants and institutional support information should be included in the thanks command. Other
% acknowledgements to colleagues or referees should appear in the Acknowledgements sections
% at the end of the paper, before the bibliography.

\author[1]{First Author\thanks{Supported by NASA grant ABC123.}}
\author[2]{Some Second Author}
\author[3]{Some Third Author}

% Write your affiation
\affil[1]{%
% Name of Department, Name of Institution, City and / or State, Country
Department of Inconsequential Studies, Solatido College, North Kentucky, U.S.A.

% leave a blank line and put your email here 
\email{fa@solatido.edu}%
}

% In case some authors share the same affiliation, use the following style.
% Do not use 'merging' of emails from the same institution such as \email{\{ssa,sta\}@uwn.edu.au}%
\affil[2,3]{%
% Name of Department, Name of Institution, City and / or State, Country
School of Hard Knocks, University of Western Nowhere, Somewhere, Australia

% leave a blank line and write the emails in the author's order
\email{ssa@uwn.edu.au}, \email{sta@uwn.edu.au}%
}

%%%%%%%%%%%%%%%%%%%%%%%%%%%%%%%%%%%%%%%%%%%%%%%%%%%
%%%%%%%%%%%%%%%%%%%%%%%%%%%%%%%%%%%%%%%%%%%%%%%%%%%

% AUTHORS COMMANDS
% Enter here all commands you need (and no more) in this paper.
% \newcommand{...}{...}

%%%%%%%%%%%%%%%%%%%%%%%%%%%%%%%%%%%%%%%%%%%%%%%%%%%
%%%%%%%%%%%%%%%%%%%%%%%%%%%%%%%%%%%%%%%%%%%%%%%%%%%

\begin{document}
%\writedatatofile

\maketitle

% ABSTRACT
% CT papers must include an abstract. The abstract should consist of a
% succinct statement of background followed by a listing of the
% principal new results that are to be found in the paper. The abstract
% should be informative, clear, and as complete as possible. Phrases
% like "we investigate..." or "we study..." should be kept to a minimum
% in favor of "we prove that..."  or "we show that...".  Do not
% include equation numbers, unexpanded citations (such as "[23]"), or
% any other references to things in the paper that are not defined in
% the abstract. The abstract may be distributed without the rest of the
% paper so it must be entirely self-contained.  Try to include all words
% and phrases that someone might search for when looking for your paper.
% You can use some basic LaTeX commands in the abstract, but not any
% user defined macros. 

\begin{abstract}
  The reconstruction conjecture states that the multiset of 
  vertex-deleted subgraphs of a graph determines the graph, provided
  it has at least 3 vertices.  This problem was independently introduced
  by Stanis\l aw Ulam (1960) and Paul Kelly (1957). In this paper,
  we prove the conjecture by elementary methods. 
  It is only necessary
  to integrate the Lenkle potential of the Broglington manifold over
  the quantum supervacillatory measure in order to reduce the set of
  possible counterexamples to a small number (less than a trillion).
  A simple computer program that implements Pipletti's classification
  theorem for torsion-free Aramaic groups with simplectic socles can
  then finish the remaining cases.
\end{abstract}

% TABLE OF CONTENTS, LIST OF FIGURES, LIST OF TABLES
% Please, do not include a table of contents, a list of figures, or a
% list of tables. They will be removed by the editors (and the command
% is actually redefined in the ct.sty file).

%%%%%%%%%%%%%%%%%%%%%%%%%%%%%%%%%%%%%%%%%%%%%%%%%%%
%%%%%%%%%%%%%%%%%%%%%%%%%%%%%%%%%%%%%%%%%%%%%%%%%%%

\section{Introduction}

The reconstruction conjecture states that the multiset of unlabeled
vertex-deleted subgraphs of a graph determines the graph, provided it
has at least three vertices.  This problem was independently introduced
by Ulam~\cite{Ulam} and Kelly~\cite{Kelly}.  The reconstruction
conjecture is widely studied
\cite{Bollobas,FGH,HHRT,KSU,WS} and is very interesting
because it is. See \cite{BH} for more about the
reconstruction conjecture.

\begin{definition} 
  A graph is \emph{fabulous} if \emph{rest of definition here}.
\end{definition}

\begin{theorem}\label{Thm:FabGraphs}
  All planar graphs are fabulous.
\end{theorem}

\begin{proof}
  Suppose on the contrary that some planar graph is not fabulous.
  Then we have a contradiction.
\end{proof}

%%%%%%%%%%%%%%%%%%%%%%%

\section{Broglington Manifolds}

This section describes background information about Broglington
Manifolds.

\begin{lemma}\label{lem:Technical}
  Broglington manifolds are abundant.
\end{lemma}

\begin{proof}
  A proof is given here.
\end{proof}

%%%%%%%%%%%%%%%%%%%%%%%

\section{Proof of Theorem~\ref{Thm:FabGraphs}}

In this section we complete the proof of Theorem~\ref{Thm:FabGraphs}.

\begin{proof}[Proof of Theorem~\ref{Thm:FabGraphs}]
Let $G$ be a graph. We have
 % The align environment for multi-line equations is defined in the amsmath package.
  \begin{align}
    |X| &= a+b+c \nonumber\\
         &= \alpha\beta\gamma.
  \end{align}
  This completes the proof of Theorem~\ref{Thm:FabGraphs}.
\end{proof}

\begin{figure}[ht]
  \centering
    % Use \includegraphics to import figures; for example 
    %      \includegraphics[scale=0.6]{filename}
  \caption{Here is an informative figure.\label{fig:InformativeFigure}}
\end{figure}

%%%%%%%%%%%%%%%%%%%%%%%%%%%%%%%%%%%%%%%%%%%%%%%%%%%
%%%%%%%%%%%%%%%%%%%%%%%%%%%%%%%%%%%%%%%%%%%%%%%%%%%

% ACKNOWLEDGEMENTS
% Include acknowledgements to colleagues and referee here.
% Funding and grant support should appear in footnotes on the front page, using the 
% thanks command in the authors command (see above).

\section*{Acknowledgements}

Thanks to Professor Qwerty for suggesting the proof of
Lemma~\ref{lem:Technical} and introducing us to the work of~\cite{Bj95}
as well as~\cite{winter}.

Some of the computation of the present paper have been done using the 
SageMath software~\cite{sage}. You can find our code over here~\cite{code}.

%%%%%%%%%%%%%%%%%%%%%%%%%%%%%%%%%%%%%%%%%%%%%%%%%%%
%%%%%%%%%%%%%%%%%%%%%%%%%%%%%%%%%%%%%%%%%%%%%%%%%%%

% BIBLIOGRAPHY
% Pease provide us with a bibtex file for your bibliography
% You can find examples in ct-sample.bib. 
% Please use the correct entrytype 
%     article: any article published in a periodical like a journal article or magazine article
%     book: a book
%     booklet: like a book but without a designated publisher
%     conference: a conference paper
%     inbook: a section or chapter in a book
%     incollection: an article in a collection
%     inproceedings: a conference paper (same as the conference entry type)
%     manual: a technical manual
%     masterthesis: a Masters thesis
%     misc: used if nothing else fits
%     phdthesis: a PhD thesis
%     proceedings: the whole conference proceedings
%     techreport: a technical report, government report or white paper
%     unpublished: a work that has not yet been officially published
%
% When it exists, please add a DOI to your reference using the doi field, this will be shown as a link in the final pdf
% Similarly for preprints, please put the arXiv reference into the eprint field 
% When necessary, you can provide a link through the URL field. Please do not give both the URL and the DOI.
%
% We encourage you to use MathSciNet https://mathscinet.ams.org/mathscinet
% or other equivalent bibliography databases to get full and correct bibliographic entries.

\bibliographystyle{alphaurl}
\bibliography{ct-sample}

\end{document}